% Artigo — Simulador Relativístico (Física 4)
% Compilar com:  tectonic artigo.tex
\documentclass[10pt,a4paper,twocolumn]{article}

\usepackage{fontspec}            % XeTeX: Unicode + acentos
\usepackage[brazil]{babel}
\usepackage{amsmath}
\usepackage{amssymb}
\usepackage{geometry}
\geometry{a4paper, top=2.4cm, bottom=2.6cm, left=1.7cm, right=1.7cm, columnsep=0.65cm}
\usepackage{graphicx}
\usepackage{booktabs}
\usepackage{array}
\usepackage{caption}
\captionsetup{labelsep=period, font=small, skip=6pt}
\usepackage{titlesec}
\titleformat{\section}{\centering\normalsize}{\thesection.}{0.55em}{\MakeUppercase}
\titleformat{\subsection}{\normalsize\itshape}{\thesubsection}{0.6em}{}
\titlespacing{\section}{0pt}{14pt plus 3pt minus 2pt}{8pt plus 2pt}
\titlespacing{\subsection}{0pt}{10pt plus 3pt minus 2pt}{5pt plus 1pt}
\usepackage{tikz}
\usetikzlibrary{arrows.meta, positioning}
\usepackage[hidelinks]{hyperref}
\urlstyle{tt}

\setlength{\parskip}{0pt}
\setlength{\parindent}{1.2em}

% Referências no estilo autor-ano com recuo deslocado
\newenvironment{referencias}
  {\section*{REFERÊNCIAS}\small
   \begin{list}{}{\setlength{\leftmargin}{1.4em}\setlength{\itemindent}{-1.4em}%
                  \setlength{\itemsep}{2pt}\setlength{\parsep}{0pt}}}
  {\end{list}}

\begin{document}

\twocolumn[{
  \begin{center}
    {\large\bfseries Simulador Relativístico: Visualização Interativa em Tempo\\[0.25em]
     Real dos Efeitos Ópticos da Relatividade Especial\,\textsuperscript{$\star$}\par}
    \vspace{1.5em}
    {\bfseries Sergio Roncato Maccari\,\textsuperscript{*}}\par
    \vspace{1.1em}
    {\small\itshape \textsuperscript{*}\,Faculdade de Engenharia de Computação, Universidade
     Tecnológica Federal do Paraná,\\ PR (e-mail: sergiomaccari@alunos.utfpr.edu.br).}\par
  \end{center}
  \vspace{0.9em}
  \hrule height 0.4pt
  \vspace{1.0em}
  {\small
  \noindent\textbf{Abstract:} Special Relativity is usually taught through isolated formulas,
  and students rarely develop a visual intuition of how the world would actually look at
  speeds close to the speed of light. This work presents an interactive simulator, running
  in real time in the web browser (TypeScript + Three.js/WebGL), that turns the canonical
  equations of Special Relativity into a navigable 3D scene: Lorentz contraction, time
  dilation, relativistic aberration (via the past light cone), the relativistic Doppler
  effect and relativistic beaming, with color computed by spectral reconstruction against
  CIE response curves. All the per-frame visible physics is evaluated per vertex on the GPU,
  in GLSL shaders, and the spectral-reconstruction integral is solved once on the CPU at
  start-up, which keeps the simulation interactive even on modest hardware. Three observation modes
  (orbital laboratory, observer's eye and first person at constant velocity) let the user
  separate what is \emph{measured} from what is \emph{seen}, which is the central conceptual
  distinction of relativistic optics.\par
  \vspace{0.6em}
  \noindent\textbf{Resumo}: A Relatividade Especial costuma ser estudada por meio de fórmulas
  isoladas, e o estudante raramente desenvolve uma intuição visual de como o mundo de fato
  pareceria a velocidades próximas à da luz. Este trabalho apresenta um simulador interativo,
  executado em tempo real no navegador (TypeScript + Three.js/WebGL), que transforma as
  equações canônicas da Relatividade Especial em uma cena 3D navegável: contração de
  Lorentz, dilatação do tempo, aberração relativística (via cone de luz passado), efeito
  Doppler relativístico e \emph{beaming}, com a cor calculada por reconstrução espectral
  com base nas curvas de resposta CIE. Toda a física visível a cada quadro é avaliada por
  vértice na GPU, em \emph{shaders} GLSL, e a integral da reconstrução espectral é resolvida
  uma única vez na CPU, na inicialização, o que mantém a simulação interativa mesmo em
  máquinas modestas. Três
  modos de observação (laboratório orbital, olho do observador e primeira pessoa a
  velocidade constante) permitem separar o que se \emph{mede} do que se \emph{vê}, distinção
  conceitual central da óptica relativística.\par
  }
  \vspace{0.8em}
  {\small\itshape
  \noindent Keywords: Special Relativity; Scientific Visualization; WebGL; Doppler Effect;
  Computer Graphics.\par
  \vspace{0.25em}
  \noindent Palavras-chave: Relatividade Especial; Visualização Científica; WebGL; Efeito
  Doppler; Computação Gráfica.\par
  }
  \vspace{0.7em}
  \hrule height 0.4pt
  \vspace{1.8em}
}]

{\let\thefootnote\relax\footnotetext{$^{\star}$ Trabalho desenvolvido na disciplina de
Física 4. A versão em inglês deste artigo está disponível no mesmo repositório
(\texttt{docs/artigo-en.pdf}).}}

\section{Introdução}

A Relatividade Especial, formulada por Einstein (1905), é apresentada nos cursos de
graduação principalmente na forma de equações: a contração do comprimento, a dilatação do
tempo e o efeito Doppler aparecem como fórmulas isoladas, verificadas em exercícios
numéricos. O estudante raramente forma uma imagem mental coerente de \emph{como o mundo
pareceria} a um viajante relativístico, em parte porque nenhuma experiência cotidiana se
aproxima do regime $v \sim c$, em parte porque os efeitos nunca ocorrem isoladamente: a
contração, o atraso de propagação da luz, a aberração, o desvio de cor e a variação de
brilho são manifestações simultâneas de uma única estrutura geométrica, o espaço-tempo de
Minkowski.

Há ainda uma confusão conceitual recorrente no estudo introdutório da Relatividade
Especial: tratar como sinônimos \emph{medir} e \emph{ver}. A contração de Lorentz é um fato
de \emph{medida} (registro simultâneo, no referencial do observador, das duas extremidades
de um objeto); já a \emph{imagem} captada por um olho ou câmera combina eventos
passados, pois a luz tem velocidade finita. Terrell (1959) e Penrose (1959) mostraram,
independentemente, que um objeto rápido de pequeno tamanho angular não parece achatado em
uma fotografia, e sim girado. Distinguir esses dois conceitos exige visualizar os
dois estágios separadamente, o que é inviável com figuras estáticas.

Atacar essa lacuna por meio de visualização interativa não é uma ideia nova, e a
referência mais influente vem do MIT Game Lab. O jogo \emph{A Slower Speed of Light}
(Kortemeyer et al., 2013) coloca o jogador em um mundo onde a velocidade da luz cai
progressivamente até a escala do caminhar, e o \emph{OpenRelativity} (MIT Game Lab, 2013;
Sherin et al., 2016) generaliza aquele motor em um conjunto de ferramentas de código
aberto para a \emph{engine} Unity. Esses projetos consolidaram o repertório visual hoje
canônico da área, isto é, contração, aberração, Doppler e \emph{beaming} aplicados a uma
cena navegável, e são a inspiração declarada deste trabalho. A diferença de recorte está
no meio de execução e no público-alvo: o OpenRelativity pressupõe a instalação da Unity e
a compilação de um projeto, ao passo que o simulador aqui descrito foi escrito do zero
como uma página \emph{web} estática, em TypeScript e GLSL sobre WebGL, que abre em
qualquer navegador moderno sem instalação e sem hardware dedicado. A prioridade de projeto
foi baixar ao máximo a barreira de acesso do estudante de graduação que quer simplesmente
ver o fenômeno.

Este trabalho apresenta um simulador interativo que preenche essa lacuna. Trata-se de uma
aplicação \emph{web} de página única, escrita em TypeScript sobre a biblioteca Three.js
(WebGL), na qual toda a física relativística é avaliada \emph{por vértice} na placa gráfica,
em \emph{shaders} GLSL. O usuário ajusta a fração da velocidade da luz,
$\beta = v/c$, e observa imediatamente o efeito combinado (ou isolado, pois cada efeito
pode ser ligado e desligado individualmente) da contração de Lorentz, do tempo retardado
(posição aparente), do efeito Doppler relativístico e do \emph{beaming} sobre uma mesma
cena tridimensional.

O recorte físico foi restrito ao domínio da Relatividade Especial em regime
simples, coerente com o escopo de Física 4: fontes e observador movem-se com
\emph{velocidade constante}, sem aceleração relativística e sem efeitos de Relatividade
Geral. Essa restrição é suficiente para reproduzir corretamente os fenômenos visuais mais
marcantes do movimento a alta velocidade, e a fundamentação teórica correspondente
(Nussenzveig, 2014; Griffiths, 2017) é desenvolvida na Seção~2. A Seção~3 descreve a
metodologia de implementação (arquitetura, algoritmos na GPU e modos de observação), a
Seção~4 apresenta os resultados e as limitações do modelo, e a Seção~5 conclui.

\section{Fundamentação Teórica}

\subsection{Postulados e transformação de Lorentz}

A Relatividade Especial repousa sobre dois postulados (Einstein, 1905): (i) as leis da
física têm a mesma forma em todos os referenciais inerciais; (ii) a luz propaga-se no vácuo
com a mesma velocidade $c$ em todos os referenciais inerciais, independentemente do
movimento da fonte ou do observador. O segundo postulado, imposto pela evidência
experimental de Michelson e Morley (1887), rompe com a cinemática galileana e obriga a
revisão dos conceitos de espaço, tempo e simultaneidade.

Para dois referenciais inerciais $S$ e $S'$, com $S'$ movendo-se com velocidade $v$ ao
longo do eixo $x$ de $S$ e origens coincidentes em $t = t' = 0$, os postulados levam à
\emph{transformação de Lorentz}:
\begin{equation}
  \begin{aligned}
    t' &= \gamma\left(t - \frac{v\,x}{c^{2}}\right), \qquad &
    x' &= \gamma\,(x - v\,t),\\
    y' &= y, & z' &= z,
  \end{aligned}
  \label{eq:lorentz}
\end{equation}
com as grandezas adimensionais centrais de toda a teoria:
\begin{equation}
  \beta = \frac{v}{c} \in [0,1), \qquad
  \gamma = \frac{1}{\sqrt{1-\beta^{2}}} \geq 1.
  \label{eq:gamma}
\end{equation}
A transformação preserva o intervalo invariante
$s^{2} = c^{2}t^{2} - x^{2} - y^{2} - z^{2}$, o que garante $c$ constante para todos os
observadores. No simulador, trabalha-se sempre com a velocidade adimensional
$\vec{\beta} = \vec{v}/c$, de modo que a geometria nunca depende explicitamente de $c$; o
valor numérico de $c$ só reaparece na conversão de intervalos para segundos.

\subsection{Contração do comprimento e dilatação do tempo}

Da Eq.~\eqref{eq:lorentz} seguem os dois efeitos cinemáticos clássicos. Uma régua de
comprimento próprio $L_{0}$, em repouso em $S'$ e alinhada ao movimento, tem as duas
extremidades registradas \emph{no mesmo instante} por um observador em $S$; o comprimento
medido é
\begin{equation}
  L = \frac{L_{0}}{\gamma} = L_{0}\sqrt{1-\beta^{2}}.
  \label{eq:contracao}
\end{equation}
Um relógio em repouso em $S'$, marcando dois eventos separados pelo tempo próprio
$\Delta\tau$, é medido em $S$ com o intervalo dilatado
\begin{equation}
  \Delta t = \gamma\,\Delta\tau \geq \Delta\tau .
  \label{eq:dilatacao}
\end{equation}
Ambos os efeitos são exibidos diretamente pelo simulador: a contração, aplicada à geometria
por vértice; a dilatação, por dois relógios analógicos no painel de leitura, com
$d\tau = dt/\gamma$ integrado a cada quadro.

\subsection{Adição relativística de velocidades}

Se uma partícula tem velocidade $\vec{u}$ (em unidades de $c$) num referencial que, por sua
vez, se move com $\vec{v}$ em relação ao laboratório, a velocidade composta não é a soma
galileana. Decompondo $\vec{u}$ nas componentes paralela e perpendicular a $\vec{v}$,
\begin{equation}
  \vec{u}\,' =
  \frac{\vec{u}_{\parallel} + \vec{v} + \vec{u}_{\perp}/\gamma_{v}}
       {1 + \vec{u}\cdot\vec{v}},
  \qquad
  \gamma_{v} = \frac{1}{\sqrt{1-|\vec{v}|^{2}}}.
  \label{eq:veladd}
\end{equation}
A Eq.~\eqref{eq:veladd} recupera o limite galileano quando $|\vec{u}|, |\vec{v}| \ll 1$ e
satura em $c$: para movimentos colineares, $u' = (u+v)/(1+uv)$, de modo que $u = 1$ (luz)
implica $u' = 1$ para qualquer $v$; o segundo postulado é respeitado automaticamente. O
fator $1/\gamma_{v}$ na componente transversal é a origem geométrica da aberração da luz.
Essa equação é implementada no núcleo de física do simulador e replicada em GLSL no
\emph{shader} de vértice, onde é avaliada vértice a vértice para compor a velocidade de
cada objeto com a do observador, garantindo $|\vec{\beta}| < 1$ em qualquer combinação.

\subsection{Medir não é ver: o cone de luz passado}
\label{sec:conedeluz}

\emph{Medir} é uma operação idealizada e simultânea no referencial do observador; é o que a
transformação de Lorentz descreve. \emph{Ver}, por outro lado, é capturar em um único
ponto, o olho, todos os raios de luz que chegam agora. Como a luz tem velocidade
finita, os fótons que atingem o olho neste instante partiram de pontos diferentes da fonte
em \emph{instantes diferentes do passado}. A imagem não é um retrato instantâneo, mas uma
costura de eventos sobre o cone de luz passado do evento de observação.

Considere o observador na origem, no instante $t = 0$, e um ponto-fonte com posição
$\vec{r}$ (em $t=0$) e velocidade constante $\vec{u}$ (em unidades de $c$). O instante
passado $\tau \leq 0$ em que esse ponto emitiu a luz que chega agora satisfaz a
\emph{condição de tempo retardado} $|\vec{r} + \vec{u}\,\tau| = -\tau$, que elevada ao
quadrado dá a equação quadrática
\begin{equation}
  \bigl(|\vec{u}|^{2}-1\bigr)\,\tau^{2}
  + 2\,(\vec{r}\cdot\vec{u})\,\tau
  + |\vec{r}|^{2} = 0 .
  \label{eq:quadratica}
\end{equation}
Das duas raízes reais, a solução física é a raiz \emph{retardada} ($\tau \leq 0$, sobre o
cone de luz passado); a outra corresponderia ao cone futuro, sem significado para uma
observação. A \emph{posição aparente} do ponto, isto é, a direção de onde a luz realmente
chega, é
\begin{equation}
  \vec{r}_{\mathrm{ap}} = \vec{r} + \vec{u}\,\tau_{\mathrm{ret}} .
  \label{eq:aparente}
\end{equation}
O simulador resolve a Eq.~\eqref{eq:quadratica} para \emph{cada vértice} da cena, a cada
quadro. A Tabela~\ref{tab:medirver} resume a distinção conceitual que organiza toda a
implementação.

\begin{table}[t]
  \centering
  \caption{Medir versus ver: a distinção central da óptica relativística e sua
           correspondência com as etapas do simulador.}
  \label{tab:medirver}
  \footnotesize
  \setlength{\tabcolsep}{4pt}
  \begin{tabular}{@{}lll@{}}
    \toprule
    Aspecto & Medir (Lorentz) & Ver (cone de luz) \\
    \midrule
    Operação    & simultânea no referencial & captura num ponto \\
    Comprimento & contração $1/\gamma$      & rotação aparente \\
    Direção     & --- (sem efeito)          & aberração \\
    Frequência  & --- (medida local)        & Doppler $D$ \\
    Brilho      & --- (sem efeito)          & \emph{beaming} $D^{3}$ \\
    No shader   & \emph{boost} do vértice   & $\tau_{\mathrm{ret}}$, $D$, cor \\
    \bottomrule
  \end{tabular}
\end{table}

\subsection{Aberração e rotação de Terrell--Penrose}

A aberração é a mudança da direção aparente de um raio de luz na troca de referencial.
Aplicando a Eq.~\eqref{eq:veladd} a um raio de luz, obtém-se a forma fechada
\begin{equation}
  \cos\theta' = \frac{\cos\theta + \beta}{1 + \beta\cos\theta},
  \label{eq:aberracao}
\end{equation}
onde $\theta$ e $\theta'$ são os ângulos entre a linha de visada e a direção do movimento
nos dois referenciais. Para $\beta \to 1$, quase toda a esfera celeste comprime-se para a
frente: o céu inteiro parece tombar na direção do movimento. No simulador, a
Eq.~\eqref{eq:aberracao} não é escrita explicitamente: ela \emph{emerge} da combinação do
\emph{boost} de Lorentz com o deslocamento para a posição aparente da
Eq.~\eqref{eq:aparente}.

A mesma geometria explica o resultado de Terrell (1959) e Penrose (1959): a luz vinda da
face mais distante de um objeto partiu mais cedo, quando o objeto estava em outra posição;
o observador enxerga simultaneamente a face frontal e uma porção da face que
geometricamente deveria estar oculta, exatamente como veria um objeto parado, porém
girado do ângulo de aberração da linha de visada; quando o objeto é visto na direção
perpendicular ao movimento, esse ângulo $\alpha$ satisfaz
$\operatorname{sen}\alpha = \beta$. A câmera não fotografa a contração;
fotografa uma rotação. Para objetos extensos e próximos somam-se deformações e encurvamento
de linhas retas, que o simulador exibe naturalmente por aplicar a
Eq.~\eqref{eq:quadratica} vértice a vértice (Figura~\ref{fig:aberracao}).

\subsection{Efeito Doppler relativístico}

O movimento também altera a frequência da luz recebida. Com $\hat{n}$ o versor que aponta
do observador para a fonte e $\vec{\beta}$ a velocidade da fonte medida no referencial do
observador, o \emph{fator Doppler} relativístico é
\begin{equation}
  D = \frac{f_{\mathrm{obs}}}{f_{\mathrm{emit}}}
    = \frac{1}{\gamma\,\bigl(1 + \vec{\beta}\cdot\hat{n}\bigr)} ,
  \label{eq:doppler}
\end{equation}
resultado da transformação de Lorentz do quadrivetor de onda: o fator $\gamma$ é a
dilatação do tempo e o termo $1 + \vec{\beta}\cdot\hat{n}$ é o atraso/avanço geométrico.
Os casos-limite são instrutivos. Na aproximação longitudinal,
$D = \sqrt{(1+\beta)/(1-\beta)} > 1$ (desvio para o azul); no afastamento,
$D = \sqrt{(1-\beta)/(1+\beta)} < 1$ (desvio para o vermelho); e na direção transversal,
$D = 1/\gamma < 1$, o chamado \emph{Doppler transversal}, efeito puramente relativístico,
sem análogo clássico e consequência direta da dilatação do tempo. O comprimento de onda
observado escala como $\lambda_{\mathrm{obs}} = \lambda_{\mathrm{emit}}/D$.

\subsection{Beaming relativístico (efeito holofote)}

A intensidade específica $I_{\nu}/\nu^{3}$ é um invariante de Lorentz; disso decorre que a
intensidade específica observada se transforma como
\begin{equation}
  I_{\nu,\mathrm{obs}}(\nu_{\mathrm{obs}}) = D^{3}\,I_{\nu,\mathrm{emit}}(\nu_{\mathrm{emit}}),
  \label{eq:beaming}
\end{equation}
com as frequências correspondentes relacionadas por $\nu_{\mathrm{obs}} = D\,\nu_{\mathrm{emit}}$.
Na origem do expoente, um fator $D$ vem do aumento da energia de cada fóton, um da taxa de
chegada dos fótons e $D^{2}$ da concentração angular imposta pela aberração, com $D^{-1}$
do alargamento da banda de frequência completando o balanço; a intensidade integrada em
frequência escala como $D^{4}$. A consequência visual é o
\emph{efeito holofote}: o que está à frente do movimento torna-se intensamente mais
brilhante e azulado; o que está atrás escurece e avermelha.

\subsection{Cor por reconstrução espectral}
\label{sec:cor}

Aplicar o desvio Doppler a uma cor RGB exige deslocar o \emph{espectro} da luz e medir
novamente a resposta do olho, e não apenas mudar o tom no círculo cromático. O simulador
modela o espectro de emissão como a soma de três gaussianas de \emph{área unitária}, uma
por primária sRGB, centradas nos comprimentos de onda dominantes dessas primárias
($611{,}4$, $549{,}1$ e $464{,}2\,\mathrm{nm}$) e com largura a meia altura de
$40\,\mathrm{nm}$, ordem de grandeza da largura espectral de um emissor de tela real.

Sob um desvio $s = \lambda_{\mathrm{obs}}/\lambda_{\mathrm{emit}} = 1/D$, a conservação do
número de fótons implica $f_{\mathrm{obs}}(\lambda) = f_{\mathrm{emit}}(\lambda/s)/s$, de
modo que uma gaussiana permanece gaussiana com centro e largura escalados por $s$:
\begin{equation}
  \mathcal{N}(\lambda;\,\mu,\,\sigma) \;\longmapsto\;
  \mathcal{N}(\lambda;\,s\mu,\,s\sigma).
  \label{eq:lobeshift}
\end{equation}
A intensidade não entra aqui; ela é o \emph{beaming} $D^{3}$, aplicado à parte. O XYZ
observado é então a integral do espectro deslocado contra as curvas de resposta do olho
$\bar{x},\bar{y},\bar{z}$ (CIE 1931, $2^{\circ}$), para as quais se adota o ajuste
multi-lóbulo de gaussianas \emph{enviesadas} de Wyman et al. (2013).

Como todo o caminho $\mathrm{RGB} \to$ espectro $\to$ desvio $\to$ CIE $\to
\mathrm{XYZ} \to \mathrm{RGB}$ é linear no RGB de repouso, ele colapsa, para cada $s$, em
uma única matriz $3\times3$:
\begin{equation}
  \mathbf{C}(s) = \mathbf{M}_{\mathrm{XYZ}\to\mathrm{RGB}}\;\mathbf{P}(s)\;
                  \bigl[\mathbf{M}_{\mathrm{XYZ}\to\mathrm{RGB}}\,\mathbf{P}(1)\bigr]^{-1},
  \label{eq:cormatriz}
\end{equation}
onde as colunas de $\mathbf{P}(s)$ são o XYZ de cada primária deslocada. O ajuste
enviesado não tem integral elementar contra uma gaussiana, e a integração é feita
numericamente na CPU, \emph{uma única vez} na inicialização, para $512$ valores de $s$
espaçados logaritmicamente em $[1/200,\,200]$. As matrizes resultantes vão para uma
textura, e o \emph{shader} faz apenas uma interpolação e um produto matriz-vetor por
vértice, sem nenhuma exponencial.

A normalização da Eq.~\eqref{eq:cormatriz} garante $\mathbf{C}(1) = \mathbf{I}$ por
construção (verificado numericamente com erro máximo de $2{,}2\times10^{-16}$): em repouso
a cor exibida é \emph{exatamente} a cor de repouso, sem necessidade de qualquer
interpolação corretiva entre a cor original e a reconstruída.

Uma consequência física do modelo: ele \emph{não reserva} energia em infravermelho ou
ultravioleta. Em $\beta$ elevado, as três gaussianas visíveis podem ser empurradas para
fora da faixa perceptível e o objeto escurece. O comportamento é coerente para um detector
de banda limitada como o olho, embora uma fonte real de espectro largo repusesse parte do
brilho com radiação vinda de fora do visível (ver Seção~\ref{sec:limitacoes}).

\section{Metodologia}

O desenvolvimento foi estruturado em três camadas: o \emph{núcleo} de física, o
\emph{motor} de simulação e a \emph{interface}, com a física visível executada
integralmente na GPU. A solução é uma aplicação \emph{web} estática: não há servidor de
física; todo o cálculo roda no navegador do usuário.

\subsection{Pilha tecnológica}

A Tabela~\ref{tab:stack} lista as tecnologias utilizadas e seus papéis.

\begin{table}[t]
  \centering
  \caption{Tecnologias utilizadas no projeto.}
  \label{tab:stack}
  \small
  \begin{tabular}{@{}>{\raggedright\arraybackslash}p{0.27\columnwidth}p{0.62\columnwidth}@{}}
    \toprule
    Tecnologia & Papel \\
    \midrule
    TypeScript & Lógica de aplicação com tipos estáticos: estado, modos,
                 controladores e interface. \\
    Three.js   & Cena 3D sobre WebGL: câmera em perspectiva, geometrias,
                 laço de renderização. \\
    GLSL       & \emph{Shaders} de vértice e fragmento escritos à mão; toda a
                 física relativística visível. \\
    Vite       & Servidor de desenvolvimento e empacotador de produção. \\
    lil-gui    & Painel de controles (velocidade, efeitos, modos). \\
    GitHub Actions e Pages & Publicação automática do site estático a cada
                 \emph{push}. \\
    \bottomrule
  \end{tabular}
\end{table}

\subsection{Arquitetura e fluxo de dados}

O código separa três responsabilidades: \texttt{core/} (física de Relatividade Especial,
com $\gamma(\beta)$ e a adição de velocidades de Minkowski, e cenas pré-definidas),
\texttt{engine/} (orquestração: laço de renderização, modos de observação, fábrica de
materiais)
e \texttt{ui/} (painel de controles e HUD de leitura). A cada quadro, a CPU calcula o vetor
$\vec{\beta}$ do observador conforme o modo ativo e o entrega à GPU por meio de
\emph{uniforms} (variáveis de entrada dos \emph{shaders}) compartilhados por
referência entre todos os materiais da cena: uma única escrita por quadro atualiza todos
os objetos simultaneamente, sem varrer a lista de materiais. A Figura~\ref{fig:pipeline}
resume o fluxo.

\begin{figure}[t]
  \centering
  \small
  \begin{tikzpicture}[
      node distance=3.5mm,
      caixa/.style={draw, rounded corners=1pt, align=center, inner sep=4pt,
                    text width=0.82\columnwidth},
      etapa/.style={draw, rounded corners=1pt, align=center, inner sep=3pt,
                    text width=0.70\columnwidth, fill=black!4},
      seta/.style={-{Stealth[length=2mm]}, thick}
    ]
    \node[caixa] (cpu) {\textbf{CPU (TypeScript)}, a cada quadro:\\[1pt]
      painel/teclado $\to$ modo $\to$ $\vec{\beta}$, posição do observador,\\
      interruptores dos efeitos $\to$ \emph{uniforms} globais};
    \node[caixa, below=of cpu] (vs) {\textbf{GPU: \emph{vertex shader}} (por vértice)};
    \node[etapa, below=2.5mm of vs] (e1) {1. \emph{boost} de Lorentz: posição no instante
      $t'=0$ do observador\\ $\Rightarrow$ contração do comprimento,
      Eq.~\eqref{eq:contracao}};
    \node[etapa, below=of e1] (e2) {2. cone de luz passado: raiz retardada da
      Eq.~\eqref{eq:quadratica}\\ $\Rightarrow$ posição aparente, aberração};
    \node[etapa, below=of e2] (e3) {3. fator Doppler $D$ [Eq.~\eqref{eq:doppler}] e
      \emph{beaming} $D^{3}$ [Eq.~\eqref{eq:beaming}]};
    \node[etapa, below=of e3] (e4) {4. cor: matriz $\mathbf{C}(s)$ da tabela espectral
      [Eq.~\eqref{eq:cormatriz}]};
    \node[caixa, below=2.5mm of e4] (fs) {\textbf{GPU: \emph{fragment shader}}
      (interpola a cor por pixel e aplica a opacidade)};
    \draw[seta] (cpu) -- (vs);
    \draw[seta] (e4) -- (fs);
  \end{tikzpicture}
  \caption{Pipeline de um quadro: a CPU prepara poucos valores por quadro e a GPU recalcula
           a posição aparente e a cor de cada vértice de cada objeto.}
  \label{fig:pipeline}
\end{figure}

\subsection{O algoritmo por vértice}

O \emph{vertex shader} executa, para cada vértice, quatro etapas em sequência.

\paragraph{Etapa 1: boost de Lorentz e contração.} A posição do vértice é transladada
para um referencial com o observador na origem. Em seguida, busca-se a posição do objeto
\emph{no instante presente do observador}, isto é, em $t' = 0$ do
referencial que se move com $\vec{\beta}$. Impondo $t' = 0$ na transformação temporal da
Eq.~\eqref{eq:lorentz} aplicada à linha de universo do vértice, obtém-se o instante de
laboratório
\begin{equation}
  T_{p} = \frac{\vec{\beta}\cdot\vec{x}_{\mathrm{lab}}}
               {1 - \vec{\beta}\cdot\vec{\beta}_{\mathrm{obj}}},
  \label{eq:tp}
\end{equation}
e o \emph{boost} espacial, decomposto nas direções paralela e perpendicular a
$\hat{\beta}$, leva à posição
\begin{equation}
  \begin{aligned}
    \vec{r}_{0} &= \vec{X} + (\gamma - 1)\,(\hat{\beta}\cdot\vec{X})\,\hat{\beta}
                   - \gamma\,\vec{\beta}\,T_{p},\\
    \vec{X} &= \vec{x}_{\mathrm{lab}} + \vec{\beta}_{\mathrm{obj}}\,T_{p}.
  \end{aligned}
  \label{eq:boost}
\end{equation}
A contração nasce do corte de simultaneidade: como $T_{p}$ varia de vértice a vértice
(relatividade da simultaneidade), a combinação do \emph{boost} com a avaliação de cada
vértice no seu próprio $T_{p}$ reduz a coordenada paralela de um objeto em repouso no
laboratório a $x_{\parallel}/\gamma$, que é exatamente a contração da
Eq.~\eqref{eq:contracao} aparecendo na tela. A velocidade do objeto vista pelo observador
é composta pela Eq.~\eqref{eq:veladd}.

\paragraph{Etapa 2: posição aparente.} Com o efeito de aberração ligado, o shader
resolve a quadrática da Eq.~\eqref{eq:quadratica} e seleciona a raiz física
$\tau \leq 0$; para $|\vec{u}| < 1$ as duas raízes têm sempre sinais opostos e a
escolha é unívoca. O vértice é deslocado para a posição aparente da
Eq.~\eqref{eq:aparente}; guardas numéricas tratam o caso quase linear
($|\vec{u}|^{2} \approx 1$, em que a quadrática degenera).

\paragraph{Etapa 3: Doppler e beaming.} Define-se a linha de visada
$\hat{n} = \vec{r}_{\mathrm{ap}}/|\vec{r}_{\mathrm{ap}}|$ e calcula-se o fator $D$ da
Eq.~\eqref{eq:doppler} com a velocidade do objeto em relação ao observador (composta na
Etapa 1) e o $\gamma$ correspondente. O brilho é multiplicado por
$D^{3}$ quando o \emph{beaming} está ativo.

\paragraph{Etapa 4: cor.} A cor de repouso passa pela reconstrução espectral da
Seção~\ref{sec:cor}, com o deslocamento $\lambda_{\mathrm{obs}} =
\lambda_{\mathrm{emit}}/D$, e o resultado é interpolado pela rasterização até o
\emph{fragment shader}, que se limita a aplicar a opacidade.

Cada efeito (contração, aberração, Doppler, \emph{beaming}) é controlado por um
interruptor próprio no painel, passado ao shader como \emph{uniform} $0/1$ (mistura linear
no caso da contração, desvio condicional nos demais). Assim, o estudante pode ligar um
efeito de cada vez e identificar a contribuição de cada termo das equações.

\subsection{Custo computacional: por vértice versus por pixel}

A reconstrução espectral é uma integral numérica sobre a faixa visível, cara demais para
ser avaliada em tempo real. A escolha de engenharia foi movê-la inteiramente para fora do
laço de renderização: como a Eq.~\eqref{eq:cormatriz} mostra que o resultado depende de um
único parâmetro escalar $s$, a integral é resolvida na CPU na inicialização e tabelada.
Resta ao \emph{shader} uma busca em textura e um produto matriz-vetor, sem exponenciais.

A segunda escolha é \emph{onde} pagar o custo residual: por \emph{pixel} (milhões de
avaliações por quadro em tela cheia) ou por \emph{vértice} (tipicamente ordens de grandeza
menos avaliações), com a placa
interpolando linearmente o resultado ao longo de cada triângulo, no estilo do sombreamento
de Gouraud. Adotou-se o cálculo por vértice, decisivo para manter interatividade em GPUs
modestas. O preço é a possibilidade de artefatos de interpolação em malhas grosseiras;
a mitigação é geométrica: as malhas da cena são subdivididas o suficiente para que a
variação de cor entre vértices vizinhos seja pequena. Complementarmente, um seletor de
qualidade ajusta a razão de pixels do renderizador (\emph{device pixel ratio}), reduzindo o
número de fragmentos rasterizados em máquinas lentas.

\subsection{Modos de observação}

Três modos, selecionáveis no painel, definem como o vetor $\vec{\beta}$ e a câmera são
construídos a cada quadro:

\begin{enumerate}
  \item \textbf{Laboratório (orbital):} câmera orbital externa; o observador conceitual
        está na origem e $\vec{\beta}$ é montado a partir de um \emph{slider} de magnitude
        e de um eixo escolhido. É o modo ideal para ver, de fora, a contração e a
        deformação aparente da cena. Atalhos fixam $\gamma = 2$, $5$ e $10$ via
        $\beta = \sqrt{1 - 1/\gamma^{2}}$.
  \item \textbf{Olho do Observador:} o observador fica parado na origem e
        \emph{$\vec{\beta}$ acompanha a direção do olhar}: virar o olho gira a direção do
        movimento e, com ela, a região azulada (à frente) e avermelhada (atrás).
  \item \textbf{Primeira pessoa:} navegação por teclado (WASD) com velocidade
        constante: a cada instante o observador ocupa um referencial inercial bem
        definido, coerente com o recorte de Física 4, sem fase de aceleração
        relativística.
\end{enumerate}

Em todos os modos a magnitude é saturada ($|\vec{\beta}| \leq 0{,}9999$) para evitar
$\gamma \to \infty$, e um painel de leitura (HUD) exibe $\beta$, $\gamma$,
$L/L_{0} = 1/\gamma$, o fator de dilatação, a taxa de quadros (FPS) e dois relógios
analógicos, um do laboratório e um do observador, integrados com $d\tau = dt/\gamma$, o
que torna a dilatação do tempo visível em tempo real. A interface oferece os idiomas
português e inglês, selecionáveis no painel de controles.

\section{Resultados}

\subsection{Efeitos observáveis}

A Figura~\ref{fig:repouso} mostra a cena de referência em repouso ($\beta = 0$): um piso em
grade, uma régua horizontal e uma galeria de sólidos coloridos. Nela, os dois relógios do
HUD avançam juntos ($\gamma = 1$).

\begin{figure}[t]
  \centering
  \includegraphics[width=\columnwidth]{figuras/lab-repouso.png}
  \caption{Cena do laboratório em repouso ($\beta = 0$, $\gamma = 1$): geometria e cores
           originais; os dois relógios do HUD avançam no mesmo ritmo.}
  \label{fig:repouso}
\end{figure}

Com $\beta = 0{,}7$ ao longo do eixo $x$ (Figura~\ref{fig:beta07}), aparecem
simultaneamente: a contração das formas na direção do movimento
($L/L_{0} = 71{,}4\%$, HUD); o desvio Doppler, com objetos azulados no sentido de
aproximação, avermelhados no oposto e o piso exibindo o gradiente contínuo de cor ao
longo da linha de visada; o \emph{beaming}, que escurece a região de recessão; e a
dilatação do tempo nos relógios ($3{,}5$~s próprios contra $4{,}9$~s de laboratório,
fator $\times 1{,}400 = \gamma$).

\begin{figure}[t]
  \centering
  \includegraphics[width=\columnwidth]{figuras/lab-beta07.png}
  \caption{Mesma cena com $\beta = 0{,}7$: contração ($L/L_{0} = 71{,}4\%$), Doppler e
           \emph{beaming}; os relógios exibem a dilatação $\gamma = 1{,}4$.}
  \label{fig:beta07}
\end{figure}

Ligando a posição aparente (Figura~\ref{fig:aberracao}), o tempo retardado da
Eq.~\eqref{eq:quadratica} desloca cada vértice individualmente: colunas retas aparecem
encurvadas e os sólidos, girados. É a manifestação, para objetos extensos, da
rotação de Terrell--Penrose, e a diferença visual direta entre \emph{o que se mede}
(Figura~\ref{fig:beta07}) e \emph{o que se vê} (Figura~\ref{fig:aberracao}).

\begin{figure}[t]
  \centering
  \includegraphics[width=\columnwidth]{figuras/lab-aberracao.png}
  \caption{$\beta = 0{,}7$ com posição aparente (cone de luz passado) ligada: colunas
           retas aparecem encurvadas e os sólidos girados (Terrell--Penrose).}
  \label{fig:aberracao}
\end{figure}

No modo Olho do Observador (Figura~\ref{fig:olho}), o observador parado na origem ``olha''
na direção do próprio movimento: mesmo a $\beta = 0{,}25$ o azulamento à frente é nítido,
com o gradiente de cor varrendo continuamente a cena conforme a linha de visada se afasta
da direção do movimento.

\begin{figure}[t]
  \centering
  \includegraphics[width=\columnwidth]{figuras/olho-beta025.png}
  \caption{Modo Olho do Observador, $\beta = 0{,}25$ na direção do olhar: desvio para o
           azul à frente, com gradiente contínuo de cor pela cena.}
  \label{fig:olho}
\end{figure}

Em $\beta = 0{,}9$ (Figura~\ref{fig:beta09}), o deslocamento espectral é tão intenso que a
luz visível de vários objetos sai da faixa perceptível e eles escurecem, consequência
prevista do modelo de cor sem reserva espectral (Seção~\ref{sec:cor}): o azul intenso
some no ultravioleta adiante e o vermelho, no infravermelho atrás.

\begin{figure}[t]
  \centering
  \includegraphics[width=\columnwidth]{figuras/lab-beta09.png}
  \caption{$\beta = 0{,}9$ ($\gamma = 2{,}294$): objetos escurecem quando sua luz visível
           é deslocada para fora da faixa perceptível, comportamento esperado do modelo
           de cor de banda limitada.}
  \label{fig:beta09}
\end{figure}

\subsection{Desempenho}

Por concentrar a física no \emph{vertex shader} e manter o \emph{fragment shader} trivial,
a simulação sustenta taxas interativas (dezenas de quadros por segundo) mesmo em
renderização por \emph{software}, sem GPU dedicada; as capturas das figuras, feitas
nessas condições, registram entre 23 e 28~FPS no HUD. Em hardware comum com aceleração
gráfica, a taxa é limitada pelo sincronismo da tela. O seletor de qualidade reduz a razão
de pixels do renderizador quando necessário, e o medidor de FPS do HUD permite ao usuário
calibrar o ajuste.

\subsection{Limitações e simplificações}
\label{sec:limitacoes}

A Tabela~\ref{tab:limitacoes} consolida as simplificações adotadas, sua natureza e o
efeito principal de cada uma. Duas delas merecem comentário. O modelo de cor não reserva
energia em infravermelho/ultravioleta: em $\beta$ alto, objetos escurecem em vez de exibir
o brilho que uma fonte real de espectro largo manteria, embora o comportamento qualitativo
do desvio (azular/avermelhar) permaneça correto. Além disso, o escopo restringe-se a
velocidade constante: as transformações aplicadas a cada quadro são rigorosamente corretas
para o referencial inercial instantâneo, mas as trocas de velocidade são idealizadas, sem
a fase de aceleração contínua e sem os fenômenos associados, como a precessão de Thomas,
restrição alinhada ao nível de Física 4.

\begin{table}[t]
  \centering
  \caption{Principais simplificações do modelo e seus efeitos.}
  \label{tab:limitacoes}
  \footnotesize
  \setlength{\tabcolsep}{3pt}
  \begin{tabular}{@{}>{\raggedright\arraybackslash}p{0.34\columnwidth}%
                     >{\raggedright\arraybackslash}p{0.19\columnwidth}%
                     >{\raggedright\arraybackslash}p{0.38\columnwidth}@{}}
    \toprule
    Simplificação & Natureza & Efeito principal \\
    \midrule
    Cor sem reserva IR/UV     & modelo de cor  & escurecimento em $\beta$ alto \\
    Sem iluminação/sombras    & renderização   & aparência autoluminosa \\
    Velocidade constante      & escopo (RE)    & sem aceleração relativística \\
    Cor por vértice (Gouraud) & desempenho     & artefatos em malha grosseira \\
    Unidades arbitrárias, $c$ ajustável & didática & só razões ($\beta$, $\gamma$, $D$) fiéis \\
    \bottomrule
  \end{tabular}
\end{table}

\subsection{Disponibilidade}

O simulador está publicado e pode ser executado em qualquer navegador moderno, sem
instalação:
\begin{center}\small
  \url{https://sergiomaccari.github.io/Simulador-Relativ-stico/}
\end{center}
O código-fonte completo (TypeScript, GLSL e este documento) está disponível em:
\begin{center}\small
  \url{https://github.com/sergiomaccari/Simulador-Relativ-stico}
\end{center}
Para execução local bastam \texttt{npm install} e \texttt{npm run dev} na raiz do
repositório.

\section{Conclusão}

O simulador desenvolvido demonstra que os efeitos ópticos da Relatividade Especial
(contração de Lorentz, dilatação do tempo, aberração via cone de luz passado, efeito
Doppler relativístico e \emph{beaming}) podem ser visualizados de forma interativa,
quantitativa e fisicamente fundamentada em uma aplicação \emph{web} aberta, sem
instalação. A implementação por vértice na GPU mostrou-se suficiente para manter
interatividade mesmo sem aceleração gráfica dedicada. Os interruptores individuais de
cada efeito, junto da distinção explícita entre \emph{medir} (transformação de Lorentz) e
\emph{ver} (tempo retardado), dão ao estudante um laboratório visual em que cada termo das
equações canônicas tem uma contrapartida observável na tela. Extensões naturais incluem a reserva espectral em
infravermelho/ultravioleta no modelo de cor, o cálculo do Doppler por fragmento e um modo
avançado com movimento acelerado (hiperbólico), este último já fora do escopo de Física 4.

\begin{referencias}
  \item Einstein, A. (1905). Zur Elektrodynamik bewegter Körper. \emph{Annalen der
        Physik}, 17, 891--921.
  \item Griffiths, D.J. (2017). \emph{Introduction to Electrodynamics}. 4th ed.
        Cambridge University Press, Cambridge.
  \item Kortemeyer, G., Tan, P. e Schirra, S. (2013). A Slower Speed of Light: developing
        intuition about special relativity with games. In \emph{Proceedings of the 8th
        International Conference on the Foundations of Digital Games (FDG 2013)},
        400--402. SASDG, Chania, Grécia.
  \item Michelson, A.A. e Morley, E.W. (1887). On the relative motion of the Earth and
        the luminiferous ether. \emph{American Journal of Science}, 34, 333--345.
  \item MIT Game Lab (2013). \emph{OpenRelativity} [software]. Massachusetts Institute of
        Technology. Licença MIT. \url{https://github.com/MITGameLab/OpenRelativity}.
  \item Nussenzveig, H.M. (2014). \emph{Curso de Física Básica, vol. 4: Ótica,
        Relatividade, Física Quântica}. 2ª ed. Blucher, São Paulo.
  \item Penrose, R. (1959). The apparent shape of a relativistically moving sphere.
        \emph{Proceedings of the Cambridge Philosophical Society}, 55(1), 137--139.
  \item Sherin, Z.W., Cheu, R., Tan, P. e Kortemeyer, G. (2016). Visualizing relativity:
        the OpenRelativity project. \emph{American Journal of Physics}, 84(5), 369--374.
  \item Terrell, J. (1959). Invisibility of the Lorentz contraction. \emph{Physical
        Review}, 116(4), 1041--1045.
  \item Wyman, C., Sloan, P.-P. e Shirley, P. (2013). Simple analytic approximations to
        the CIE XYZ color matching functions. \emph{Journal of Computer Graphics
        Techniques}, 2(2), 1--11.
\end{referencias}

\end{document}
